%% toc: P2 10 10.1,10.2
\section{Proofs for congruence and replacement}
\label{sec:annex-p02-w02-p02-003-exercise-link}

Notation is that of \cref{def:p02-selected-replacement}.  For a formula
\(\varphi\) and a set \(O\) of positions, \(O\) is \emph{admissible} for
\(\varphi\) and \(\psi\) when
\(O\subseteq\{a\in\operatorname{Pos}(\varphi):\varphi_a=\psi\}\).

\begin{annexitem}{annex:p02-w02-p02-003-replace}{thm:p02-w02-p02-003-congruence-replacement}{Congruence and replacement}
\begin{proof}
Two facts about the replacement recursion are used in both directions.

\emph{Fact 1: the empty choice leaves a formula unchanged.}  For every
formula \(\varphi\) and every \(\chi\), \(\varphi[\varnothing\leftarrow\chi]=\varphi\).
We argue by structural induction on \(\varphi\).  The root is not in
\(\varnothing\), and both sets \(O_1,O_2\) computed from \(\varnothing\) are
empty.  For \(p_n\) and \(\bottom\) the first two recursion clauses return the
leaf itself.  For \(\neg\alpha\) the recursion returns
\(\neg(\alpha[\varnothing\leftarrow\chi])\), which is \(\neg\alpha\) by the
induction hypothesis.  For \((\alpha\circ\beta)\) with
\(\circ\in\{\land,\lor,\to\}\) it returns
\((\alpha[\varnothing\leftarrow\chi]\circ\beta[\varnothing\leftarrow\chi])\),
which is \((\alpha\circ\beta)\) by the two induction hypotheses.

\emph{Fact 2: admissibility passes to immediate subformulas.}  Suppose \(O\)
is admissible for \(\varphi\) and \(\psi\), and \(\epsilon\notin O\).  If
\(\varphi=\neg\alpha\), let \(a\in O_1\).  Then \(1a\in O\), so
\(1a\in\operatorname{Pos}(\neg\alpha)\), which by the definition of
\(\operatorname{Pos}\) means \(a\in\operatorname{Pos}(\alpha)\).  Also
\(\alpha_a=(\neg\alpha)_{1a}=\psi\).  Hence \(O_1\) is admissible for
\(\alpha\) and \(\psi\).  If \(\varphi=(\alpha\circ\beta)\), the same argument
with \((\alpha\circ\beta)_{1a}=\alpha_a\) and
\((\alpha\circ\beta)_{2a}=\beta_a\) shows that \(O_1\) is admissible for
\(\alpha\) and \(\psi\), and \(O_2\) for \(\beta\) and \(\psi\).

\emph{(a) implies (b).}  Assume (a), and fix \(\psi,\chi\) with
\(\psi\mathrel{E}\chi\).  We prove by structural induction on \(\varphi\)
that \(\varphi\mathrel{E}\varphi[O\leftarrow\chi]\) for every admissible
\(O\).  The induction hypothesis for a proper subformula covers every
admissible set for that subformula.

If \(\epsilon\in O\), then \(\varphi=\varphi_\epsilon=\psi\) by
admissibility, and the recursion returns \(\chi\).  The required statement is
\(\psi\mathrel{E}\chi\), which is the hypothesis.

Assume now that \(\epsilon\notin O\).
\begin{itemize}
  \item \(\varphi=p_n\) or \(\varphi=\bottom\).  The only position is
        \(\epsilon\), so \(O=\varnothing\).  The recursion returns
        \(\varphi\), and reflexivity of \(E\) gives
        \(\varphi\mathrel{E}\varphi\).
  \item \(\varphi=\neg\alpha\).  By Fact~2, \(O_1\) is admissible for
        \(\alpha\), so the induction hypothesis gives
        \(\alpha\mathrel{E}\alpha[O_1\leftarrow\chi]\).  Negation congruence
        in (a) gives
        \(\neg\alpha\mathrel{E}\neg(\alpha[O_1\leftarrow\chi])\), and the
        right-hand formula is \(\varphi[O\leftarrow\chi]\).
  \item \(\varphi=(\alpha\circ\beta)\) with \(\circ\in\{\land,\lor,\to\}\).
        By Fact~2 and the induction hypothesis,
        \(\alpha\mathrel{E}\alpha[O_1\leftarrow\chi]\) and
        \(\beta\mathrel{E}\beta[O_2\leftarrow\chi]\).  The congruence clause
        of (a) for this \(\circ\) gives
        \((\alpha\circ\beta)\mathrel{E}
         (\alpha[O_1\leftarrow\chi]\circ\beta[O_2\leftarrow\chi])\), and the
        right-hand formula is \(\varphi[O\leftarrow\chi]\).
\end{itemize}
These are all the formation clauses of \(\mathsf{Fm}_{\mathsf{CPL}}\), so (b)
holds.  This direction uses reflexivity of \(E\) and the four congruence
clauses for \(\neg,\land,\lor,\to\).  It uses neither symmetry nor
transitivity.

\emph{(b) implies (a).}  Assume (b).  The clause
\(\bottom\mathrel{E}\bottom\) is reflexivity of \(E\).

For negation, let \(\alpha\mathrel{E}\alpha'\).  In \(\varphi=\neg\alpha\)
the position \(1\) exists, since \(\epsilon\in\operatorname{Pos}(\alpha)\),
and \(\varphi_1=\alpha_\epsilon=\alpha\).  So \(O=\{1\}\) is admissible for
\(\varphi\) and \(\psi=\alpha\).  With \(\chi=\alpha'\), the root of
\(\varphi\) is unselected and \(O_1=\{\epsilon\}\).  The recursion gives
\(\varphi[O\leftarrow\alpha']=\neg(\alpha[\{\epsilon\}\leftarrow\alpha'])
=\neg\alpha'\), because a selected root returns the replacing formula.  By
(b), \(\neg\alpha\mathrel{E}\neg\alpha'\).

For a binary connective \(\circ\in\{\land,\lor,\to\}\), let
\(\alpha\mathrel{E}\alpha'\) and \(\beta\mathrel{E}\beta'\).
\begin{itemize}
  \item In \((\alpha\circ\beta)\) the set \(\{1\}\) is admissible for the
        formula \(\alpha\), because \((\alpha\circ\beta)_1=\alpha\).
        Replacing by \(\alpha'\), the root is unselected, \(O_1=\{\epsilon\}\)
        and \(O_2=\varnothing\).  The recursion and Fact~1 give
        \((\alpha\circ\beta)[\{1\}\leftarrow\alpha']
         =(\alpha'\circ\beta)\).  By (b),
        \((\alpha\circ\beta)\mathrel{E}(\alpha'\circ\beta)\).
  \item In \((\alpha'\circ\beta)\) the set \(\{2\}\) is admissible for the
        formula \(\beta\), because \((\alpha'\circ\beta)_2=\beta\).
        Replacing by \(\beta'\), the root is unselected, \(O_1=\varnothing\)
        and \(O_2=\{\epsilon\}\).  Fact~1 and the recursion give
        \((\alpha'\circ\beta)[\{2\}\leftarrow\beta']
         =(\alpha'\circ\beta')\).  By (b),
        \((\alpha'\circ\beta)\mathrel{E}(\alpha'\circ\beta')\).
\end{itemize}
Transitivity of \(E\) gives
\((\alpha\circ\beta)\mathrel{E}(\alpha'\circ\beta')\).  All primitive
connectives have now been treated, so (a) holds.  This direction uses
reflexivity and transitivity of \(E\) and does not use symmetry.
\end{proof}
\end{annexitem}
